% !TeX root = ../../relatorio-projeto-exemplo-undf.tex
\undfsetcontent{resumo}{%
Este relatório didático apresenta o desenvolvimento de modelos de relatório
final no projeto UnDF-TeX. O problema identificado foi a presença de um roteiro
de planejamento em uma modalidade cuja entrega deveria relatar trabalho
realizado, resultados e limitações. O objetivo foi organizar modelos para
projeto técnico ou profissional e para análise de negócio tecnológico, com
exemplos independentes. O procedimento compreendeu inspeção dos arquivos,
reescrita dos roteiros, separação entre conteúdo editável e apresentação
institucional e compilação com LuaLaTeX e Biber. A avaliação documental
considerou quatro entradas compiláveis, presença dos elementos institucionais,
resumos em português e inglês, sumário e referências. Foram produzidos dois
modelos para preenchimento e dois exemplos, com sete seções primárias no
projeto e oito no negócio tecnológico. A apresentação foi compartilhada,
enquanto metadados, texto e bibliografia permaneceram independentes. O
resultado atende à organização técnica proposta, sem demonstrar ganhos de
usabilidade ou eficácia acadêmica. Não houve pesquisa com estudantes,
implantação institucional ou avaliação por banca. Concluiu-se que os artefatos
permitem relatar entregas finais com limites explícitos, permanecendo necessária
a conferência das exigências específicas de cada curso antes do uso acadêmico.}
\undfsetcontent{palavraschave}{relatório final; modelos acadêmicos; LaTeX}
\undfsetcontent{abstract}{%
This educational report describes the development of final report templates
within the UnDF-TeX project. The identified problem was a planning outline
being used for a modality that required an account of completed work,
findings and limitations. The objective was to organize templates for a
technical or professional project and for a technology business analysis,
with independent examples. The procedure included inspecting existing files,
rewriting the outlines, separating editable content from institutional
presentation and compiling the documents with LuaLaTeX and Biber. Document
evaluation considered four compilable entry points, institutional elements,
Portuguese and English abstracts, a table of contents and references. Two
editable templates and two examples were produced, with seven primary sections
for the project and eight for the technology business. Presentation commands
were shared, while metadata, text and bibliographic records remained
independent. The result meets the proposed technical organization without
demonstrating usability gains or academic effectiveness. No student research,
institutional deployment or examination board assessment was conducted. The
artifacts provide a structure for describing completed work with explicit
limitations. Requirements specific to each degree program still need to be
checked before academic use.}
\undfsetcontent{keywords}{final report; academic templates; LaTeX}
